\documentclass[UTF8,zihao=-4,a4paper,oneside,openany,fontset=none]{ctexbook}
\usepackage[left=24mm,right=23mm,top=24mm,bottom=24mm,headheight=15pt]{geometry}
\usepackage{fontspec}
\setmainfont{Liberation Serif}
\setsansfont{Arimo}
\setmonofont{DejaVu Sans Mono}[Scale=.82]
\setCJKmainfont{Noto Serif CJK SC}[ItalicFont={Noto Serif CJK SC}]
\setCJKsansfont{Noto Sans CJK SC}
\setCJKmonofont{Noto Sans Mono CJK SC}
\newCJKfontfamily\heiti{Noto Sans CJK SC}
\IfFontExistsTF{Arial}{\newfontfamily\figfont{Arial}}{\newfontfamily\figfont{Arimo}}
\usepackage{amsmath,amssymb,amsthm,mathtools,bm}
\usepackage{graphicx,booktabs,longtable,array,tabularx,multirow}
\usepackage{float,placeins,caption,subcaption}
\usepackage{tikz}
\usetikzlibrary{arrows.meta,positioning,shapes.geometric,fit}
\usepackage{listings,xcolor}
\usepackage[numbers,sort&compress]{natbib}
\usepackage[hidelinks,unicode]{hyperref}
\usepackage{xurl,fancyhdr,enumitem}
\usepackage[ruled,vlined,linesnumbered]{algorithm2e}
\usepackage{microtype}
\hypersetup{pdftitle={质谱与DEMS信号解析的数学方法及智能体接口},pdfauthor={Convolvex-MS},pdfsubject={MCP, inverse problems, mass spectrometry}}
\ctexset{chapter={format=\centering\heiti\Large,beforeskip=12pt,afterskip=22pt},section={format=\heiti\large},subsection={format=\heiti\normalsize}}
\linespread{1.12}
\setlength{\parskip}{.10em}
\setlength{\parindent}{2em}
\setlength{\emergencystretch}{2em}
\pagestyle{fancy}\fancyhf{}\fancyhead[L]{\small 质谱与 DEMS 信号解析}\fancyhead[R]{\small Convolvex--MS 1.1}\fancyfoot[C]{\thepage}
\renewcommand{\headrulewidth}{.35pt}
\captionsetup{font=small,labelfont=bf,labelsep=quad}
\numberwithin{equation}{chapter}
\theoremstyle{definition}\newtheorem{definition}{定义}[chapter]\newtheorem{example}{例}[chapter]
\theoremstyle{definition}\newtheorem{proposition}{命题}[chapter]
\renewcommand{\proofname}{证明}
\newcommand{\R}{\mathbb{R}}\newcommand{\Aop}{\mathcal{A}}\newcommand{\prox}{\operatorname{prox}}
\newcommand{\diag}{\operatorname{diag}}\newcommand{\argmin}{\operatorname*{arg\,min}}
\newcommand{\code}[1]{\texttt{\detokenize{#1}}}
\newcommand{\fig}[3]{\begin{figure}[htbp]\centering\includegraphics[width=#2\linewidth]{figures/#1.pdf}\caption{#3}\label{fig:#1}\end{figure}}
\lstset{basicstyle=\ttfamily\small,breaklines=true,columns=fullflexible,keepspaces=true,showstringspaces=false,frame=single,rulecolor=\color{black!25},framesep=5pt,aboveskip=10pt,belowskip=10pt}
\SetAlgorithmName{算法}{算法}{算法目录}\SetKwInput{KwIn}{输入}\SetKwInput{KwOut}{输出}
\graphicspath{{figures/}}
\newcommand{\TestPassed}{147}
\newcommand{\TestSkipped}{2}

\begin{document}
\hypersetup{pageanchor=false}
\begin{titlepage}\centering
\vspace*{25mm}
{\heiti\LARGE 质谱与 DEMS 信号解析\par}
\vspace{9mm}{\heiti\LARGE 的数学方法及智能体接口\par}
\vspace{19mm}{\Large Convolvex--MS 1.1.0\par}
\vspace{10mm}{\large 技术报告\par}
\vfill
\begin{minipage}{.78\textwidth}\normalsize
从观测方程、传输过程与离散算子出发，建立质谱分解、正则化反卷积、背景消元与定量审查的统一方法，并给出 MCP 服务及可审计智能体工作流的实现与验收记录。
\end{minipage}
\vfill{\large 2026 年 11 月 14 日\par}\vspace{12mm}
\end{titlepage}
\frontmatter
\hypersetup{pageanchor=true}
\chapter*{编写说明}\addcontentsline{toc}{chapter}{编写说明}
本报告以实际交付的 Convolvex--MS 1.1.0 源码为对象。讨论范围限于线性质谱混合、质谱质量轴线形、DEMS 时间响应、背景与同位素诊断、定量条件及智能体调用。通用图像去模糊、三维图像恢复与前端设计均不列入本报告。

二维数组不等于二维图像。一个以质荷比为行、以时间为列的观测矩阵正是质谱数据的常见表示；本报告保留此类矩阵和热图，但不沿离散质量通道任意施加图像平滑。通用数学内核中的兼容代码仍可保留，其存在不构成质谱分析需要图像恢复的论据。

正文按观测量、物理模型、离散化、反问题、数值算法、定量审查、接口、工作流和实验验收的顺序展开。命题中的条件与结论均有明确边界。程序计算出的“收敛解”、化学物种的“可识别性”和实验标定后的“绝对量”分别对待。

本版全部演示与基准使用公开构造的合成数据，不使用用户原始实验数据，也不把构造的响应矩阵称为 NIST 标准谱。测试通过不能替代标准品回收、空白和重复实验。关于算法优化，分通道投影、批量 FFT 与尺度相容停止判据属于本项目的实现改进；变量投影、FISTA 和原始--对偶方法本身不是本项目首创。

绘图配置的默认字体为 Arial。当前构建环境未安装 Arial，故本次 PDF 图中文字明确采用 Arimo；该事实记录于 \code{simulations/font_receipt.json}。默认严格绘图模式遇到缺少 Arial 会报错，不会静默替换。安装本机合法字体后可由相同脚本重新生成精确 Arial 版本；交付包不附带任何字体文件。

本版实现 MCP 的 \code{2026-07-28} 核心调用路径，以及 \code{2025-11-25}、\code{2025-06-18} 兼容路径。协议测试覆盖实际本地子进程与 HTTP 测试客户端；官方 SDK 与第三方宿主尚未在本环境验证。本报告不宣称获得协议认证。
\tableofcontents
\chapter*{符号与量纲}\addcontentsline{toc}{chapter}{符号与量纲}
\begin{longtable}{p{.16\textwidth}p{.52\textwidth}p{.22\textwidth}}
\toprule 符号 & 含义 & 单位或维数\\\midrule\endhead
$m,j,n$ & 质量通道、候选组分、时间样点的指标 & 整数\\
$M,J,N$ & 通道数、组分数、样点数 & 正整数\\
$t,\Delta t$ & 物理时间、均匀时间步长 & s\\
$y_m(t)$ & 质谱观测信号 & A、counts 或声明单位\\
$r_j(t)$ & 电极侧组分生成速率 & mol\,s$^{-1}$\\
$\kappa_j$ & 包括收集效率的有效灵敏度 & C\,mol$^{-1}$\\
$x_j(t)=\kappa_jr_j(t)$ & 参考通道等效源信号 & A\\
$h_j(t)$ & 单位积分的因果传输响应密度 & s$^{-1}$\\
$h_{j,k}^{\mathrm d}$ & 离散 FIR 核权重 & 无量纲\\
$S_{mj}$ & 相对碎片响应；参考通道系数为 1 & 无量纲\\
$Y,X$ & 观测、源信号矩阵 & $M\times N$、$J\times N$\\
$B,\beta_m$ & 时间背景基、通道背景系数 & $N\times p$、$p$\\
$W_m$ & 通道内的逆标准差与掩码算子 & 逆观测单位\\
$D_t$ & 离散时间差分 & 需区分索引差分与物理导数\\
$\lambda_1,\lambda_2$ & L1 与 L2 目标函数系数 & 随量纲定义而定\\
$F$ & 法拉第常数 & $96485.33212331001$ C\,mol$^{-1}$\\
$Q^-,Q^+$ & 阴极、阳极电荷的正值大小 & C\\
\bottomrule
\end{longtable}
\mainmatter
\chapter{观测对象与问题的确定}
\section{先确定观测量，再选择算法}
质谱仪输出的是检测器的响应，不是未经变换的物质生成速率。对同一组数据，可以提出质量峰恢复、混合谱分离、时间响应反演和绝对定量四种不同问题。它们的未知量、先验与验收条件不同，不能以同一个“反卷积”名称掩盖这些差别。

\begin{definition}[质谱观测矩阵]
取质量通道 $m_1,\ldots,m_M$ 和时间样点 $t_1,\ldots,t_N$，定义
\begin{equation}
Y_{mn}=y_{m_m}(t_n),\qquad Y\in\R^{M\times N}.
\end{equation}
行代表质量通道，列代表采样时间。每一行是一条离子流曲线，每一列是一个时刻的采样谱。
\end{definition}
该矩阵可以画为热图，也可逐行画曲线。图~\ref{fig:01_mass_time} 是本报告使用的合成观测，其二维性来自两种实验坐标，而不是来自图像的两个空间坐标。质谱数据还可能具有保留时间、离子迁移时间等维度；多维数据组织并不意味着应采用空间图像模型\cite{deimos}。
\fig{01_mass_time}{.93}{合成的质量通道--时间观测矩阵。质量通道为离散类别，图中相邻两行不必对应相邻整数质量，更不表示两个空间像素。}

\section{三种线性变换不能混为一谈}
质量轴线形使一个窄质量峰展宽，数学上可在局部近似为沿质量轴的卷积。碎裂和不同分子的重叠谱构成通道混合，通常是矩阵乘法。膜、管路和采样腔的传输过程则使生成速率延迟并展宽，构成时间卷积。三者的结构为
\begin{equation}
\text{源项}\ \xrightarrow{\text{时间传输}}\ \text{组分响应}
\ \xrightarrow{\text{谱混合}}\ \text{质量通道信号},
\end{equation}
而质量轴展宽是检测器谱分辨率所对应的另一层变换。
本软件因此提供 \code{spectral}、\code{spectrotemporal} 和独立的质量轴 \code{mass-profile} 路径。第一条路径的估计量是经过传输后的组分响应，只有第二条路径在已知传输核条件下才尝试恢复源项。

\section{数据合同}
一个数值矩阵只有在坐标、单位、有效性掩码与实验条件明确后才成为可解释的观测。\code{MSDataset} 强制保存以下信息：时间数组 \code{time_s}、有序质量通道 \code{mz}、以通道为行的 \code{signal}、\code{signal_unit} 和 \code{time_origin}。电化学电流使用自己的 \code{ec_time_s} 与 \code{current_A}，不能为了方便而先插值到质谱时间轴。

同一时间轴必须严格递增；重复时间、无穷大、NaN、维度不匹配与重复质量通道均拒绝导入。非均匀时间不是错误数据，但它不能直接进入本版均匀 FFT 的时间反卷积路径。每个通道可以提供有效性掩码；掩码位置不参加拟合，也不作为零值观测。至少应有一个有效样点。

负离子流或负背景残差不在导入时删除。它可能来自基线定义、噪声、仪器偏移或真实的信号减小。非负约束只能施加于有相应物理意义的未知源项，不能据此把原始负数改为零。

\begin{example}[单位错误对正则化的影响]
同一电流用 A 与 nA 表示，数值相差 $10^9$。未加权残差平方相差 $10^{18}$；保持同一个数值的 L1 系数并不能表示同一个先验。若同时把信号和标准差转换为相同单位，白化残差才保持不变。即使如此，源项先验及核的单位仍需相应转换。
\end{example}

\section{时间同步与顺序扫描}
绝对时刻、相对起点和仪器内部时钟偏移是三个不同概念。\code{time_origin} 是数据生产者对时间坐标的明确声明，软件不会把相关峰最高的位置自动当作零时差。采样传输滞后属于物理模型，也不能一律当作时钟误差。

顺序扫描的各质量通道可能在一帧内不同时间采样。严格模型应写为
\begin{equation}
Y_{mn}=\sum_jS_{mj}(h_j*x_j)(t_n+\delta_m)+b_m(t_n+\delta_m)+\varepsilon_{mn}.
\end{equation}
当前 MCP 路径仅接受共同时间网格；非零 \code{channel_offsets_s} 会被明确拒绝，而不是忽略。一般采样算子可在 Python 层组合实现，但本版没有把任意时钟校正包装为已验证的自动功能。

\section{计算结论的层级}
本报告区分三层结论。第一层是数值结论：算法是否到达停止条件，前向重构与观测是否相符。第二层是模型结论：给定候选组分、背景和传输核时，解是否可辨识，残差是否包含结构。第三层是实验结论：候选分子是否得到独立证据支持，以及标定能否把响应转换为绝对物质量。

前一层成立不能推出后一层成立。一个错误的候选谱也可能拟合有限通道；一个条件良好的响应矩阵也不能证明样品中不存在未列入的组分。程序因此保存输入标签但不把标签当作化学鉴定结果；在结果中分别返回数值状态、科学解释状态与定量状态。

\chapter{谱混合、质量轴线形与同位素贡献}
\section{碎片响应矩阵}
电子轰击等电离过程可以使一个中性分子形成多个离子碎片。有限监测通道中的单组分指纹记为 $s_j\in\R^M$。在线性响应区间，混合物满足
\begin{equation}
 y_m(t)=\sum_{j=1}^J S_{mj}z_j(t)+b_m(t)+\varepsilon_m(t),
 \qquad Y=SZ+B_0+E.\label{eq:mix}
\end{equation}
$S$ 的第 $j$ 列是组分 $j$ 在本机条件下的相对通道响应。若选择参考通道 $m_j$，可令 $S_{m_jj}=1$，把绝对响应系数留在 $z_j$ 或后述 $\kappa_j$ 中。不同列的归一化方式必须一致地记录，不能在拟合后无痕改变。

NIST Chemistry WebBook 的谱可作为候选结构与碎片模式的参考\cite{nistweb}。相对指纹不等于当前离子源能量、质量透过率、电子倍增器增益和采样条件下的绝对灵敏度。数据库谱可以帮助定义假设，却不能替代本机标准品与混合物验证。
\fig{02_response}{.79}{本报告构造的七通道、三组分响应矩阵。A、B、C 是数学组分标签，不表示已鉴定的化合物；系数不是 NIST 谱库数据。}

\section{质量轴上的卷积}
设理想质量分布为 $a(\mu)$，仪器线形为 $g(m,\mu)$，则观测为
\begin{equation}
 p(m)=\int g(m,\mu)a(\mu)\,\mathrm d\mu+\varepsilon(m).
\end{equation}
只有在所考察窄窗口内线形近似只依赖差值时，才能写成
\begin{equation}
 g(m,\mu)=g(m-\mu),\qquad p=g*a+\varepsilon.\label{eq:massconv}
\end{equation}
宽质量范围内的分辨率变化、峰形不对称变化、质量校准漂移都可能使此近似失效。软件的质量轴路径要求均匀 \code{mz_grid}、奇数长度的居中非负线形以及单位离散和；核中心与质量格点一一对应，避免半格偏移。

式~\eqref{eq:massconv} 的恢复对象仍是一个分辨率受限的质量分布。它不自动完成多电荷包络反演、同位素簇分组、串联质谱库检索或分子结构鉴定。那些任务需要额外离子化与化学模型，未在本版 MCP 接口中实现。

\fig{07_mass_profile}{.93}{窄质量窗口中的一维线形恢复。观测与真值、恢复曲线均来自合成数据；所用线形在此窗口内固定。图像去模糊不属于本例。}

\section{以二氧化碳为例推导同位素比}
设 $p_{12},p_{13}$ 是碳同位素丰度，$p_{16},p_{17},p_{18}$ 是氧同位素丰度。只考虑单电荷分子离子，且暂不考虑质量歧视，则以 $^{12}\mathrm C{}^{16}\mathrm O_2$ 为基准，质量 45 的主要分子组合为
\begin{equation}
^{13}\mathrm C{}^{16}\mathrm O_2,
\qquad ^{12}\mathrm C{}^{16}\mathrm O{}^{17}\mathrm O.
\end{equation}
两个氧位点给出组合因子 2。因此
\begin{equation}
\frac{I_{45}}{I_{44}}
=\frac{p_{13}}{p_{12}}+2\frac{p_{17}}{p_{16}}.
\end{equation}
质量 46 的主要组合为一个 $^{18}\mathrm O$、两个 $^{17}\mathrm O$，以及 $^{13}\mathrm C$ 与一个 $^{17}\mathrm O$，故
\begin{equation}
\frac{I_{46}}{I_{44}}
=2\frac{p_{18}}{p_{16}}
+\left(\frac{p_{17}}{p_{16}}\right)^2
+2\frac{p_{13}}{p_{12}}\frac{p_{17}}{p_{16}}.\label{eq:isotope46}
\end{equation}
采用 NIST 表中代表性组成，可得到约 $1.158\%$ 和 $0.412\%$\cite{nistC,nistO}。天然材料同位素组成存在变化范围\cite{ciaaw}；仪器质量歧视、碎片、其它分子的重叠和背景也会改变实测比例。因此这些数值只是物理检查先验，不是可以无条件套用的背景扣除系数。

\section{同位素诊断与产物鉴定的区别}
在 CO$_2$ 占优的背景中，45 和 46 通道与 44 共同变化是合理的先验现象。若先在可信空白段建立参考通道、压力和缓慢漂移模型，再在反应段检查残差，可以判断是否存在超出该背景模型的附加信号。但附加信号仍可能来自未建模的系统变化。

因此，“45 有残差”与“产生了乙醇”不是等价命题。判断候选产物还需要独立碎片、标准品、空白、标记实验或其它分析方法相互支持。本版的 \code{ms_background} 输出带符号残差、回归条件和子数据集，不输出物种确认标签。

\section{谱可辨识性的初步检查}
若 $\operatorname{rank}(S)<J$，仅利用瞬时谱无法唯一分离 $J$ 个组分；即使满秩，几乎共线的列也会把小噪声放大。为消除列幅值对诊断的干扰，先令
\begin{equation}
\widetilde s_j=s_j/\|s_j\|_2,\qquad
\mu=\max_{i\ne j}|\widetilde s_i^T\widetilde s_j|.
\end{equation}
软件报告归一化奇异值、条件数、最大相干性和近共线列对。该诊断只考察 $S$；它不是完整谱--时间联合算子的可辨识性定理。不同且已知的时间核有时能补充信息，未知核与背景也可能进一步削弱信息。

\chapter{DEMS 的传输物理与响应函数}
\section{从守恒方程到检测信号}
DEMS 中产物从电极附近到达质谱检测器通常经过液相扩散、膜交换以及气相输运。以浓度场 $c(\bm r,t)$ 表示传输物种，在简化的稀溶液、定常流场条件下可写
\begin{equation}
\frac{\partial c}{\partial t}+\bm v\cdot\nabla c
=\nabla\cdot(D\nabla c)+q(\bm r,t).
\end{equation}
电极源项由生成通量给定；膜面可以用线性传质边界近似，例如
\begin{equation}
-D\frac{\partial c}{\partial n}=k_m(c-c_{\mathrm{eq}}).
\end{equation}
仪器响应是穿过采样界面的通量经后续腔体、电离和检测变换后的结果。几何、扩散系数、流速、膜状态和收集效率都进入映射关系。原始研究已明确指出，传输会改变电化学质谱的动态响应，不能把检测信号直接等同于界面反应速率\cite{krempl2021}。

\section{线性时不变近似}
\begin{definition}[因果响应核]
当源项加倍使输出加倍、多个源项的输出可相加，且系统参数不随实验时间改变时，以单位脉冲源项得到的响应定义为 $h(t)$。因果性要求 $h(t)=0$ 对一切 $t<0$ 成立。
\end{definition}
将一般源项看作脉冲叠加，有
\begin{equation}
(h*r)(t)=\int_0^{\infty}h(\tau)r(t-\tau)\,\mathrm d\tau.
\end{equation}
为分离形状和收集幅值，本报告约定 $\int h(t)\,\mathrm dt=1$，并把收集效率和绝对检测响应合并为 $\kappa$。于是
\begin{equation}
 y_m(t)=b_m(t)+\sum_j S_{mj}\kappa_j(h_j*r_j)(t)+\varepsilon_m(t),\label{eq:demsfull}
\end{equation}
令 $x_j=\kappa_jr_j$，则进入反问题的未知量为等效源信号 $x_j$。其中 $[r]=\mathrm{mol\,s^{-1}}$、$[h]=\mathrm{s^{-1}}$、$[\kappa]=\mathrm{C\,mol^{-1}}$，故 $[x]=[y]=\mathrm A$，量纲闭合。

程序允许未归一化核参与一般拟合，并记录其质量；它不会暗中改变核。只有进入绝对定量路径时，才要求核质量、相对响应归一化和灵敏度定义相互一致。

\section{两个可计算的响应模型}
一阶充分混合模型满足
\begin{equation}
\tau\frac{\mathrm dz}{\mathrm dt}+z=x,
\qquad h(t)=\frac1\tau e^{-t/\tau}\mathbf1_{t\ge0}.
\end{equation}
该模型描述一个时间尺度，适合教学和某些集总近似，但并不等同于全部膜传输过程。

另考虑一维恒定漂移 $v>0$、扩散系数 $D>0$，从起点到距离 $L>0$ 处吸收边界的首次到达问题。其通量响应为
\begin{equation}
 h(t)=\frac{L}{\sqrt{4\pi Dt^3}}
 \exp\!\left[-\frac{(L-vt)^2}{4Dt}\right],\qquad t>0.
\end{equation}
它的平均到达时间为 $L/v$，方差为 $2DL/v^3$。这个结果由特定边界问题推得，说明平移与展宽可以来自同一传输过程；它不应在未检查几何时直接套用于任意 DEMS 池。
\fig{03_kernels}{.89}{联合模拟采用的三个离散响应核。它们由不同时间尺度的指数形状截断并按离散和归一化，均为合成设定。}

\section{带宽、噪声与有限时间窗}
傅里叶变换使线性卷积变为
\begin{equation}
\widehat y(\omega)=\widehat h(\omega)\widehat x(\omega)+\widehat\varepsilon(\omega).
\end{equation}
直接相除得到的误差为 $\widehat\varepsilon/\widehat h$。在 $|\widehat h|$ 很小的频率，高频噪声被急剧放大。正则化的作用是在数据支持与先验之间取得稳定折中，而不是凭空增加仪器带宽。
\fig{11_transfer_bandwidth}{.89}{合成传输核的频率响应。随频率增加而衰减的幅值意味着快速变化信息更难稳定恢复；图中的衰减不是独立噪声估计。}

若观测在响应尾部结束前停止，$\int_{t_0}^{t_1}(h*x)\,\mathrm dt$ 一般小于完整响应积分。反卷积也不能唯一确定没有观测到的任意后续源项。软件的定量路径要求操作者确认响应窗完整，但该布尔声明不构成物理完整性的证明，仍需检查尾部与实验记录。

\section{模型失效的情形}
流量或总压改变导致收集效率变化、膜润湿发生变化、信号进入饱和区、离子源条件改变，以及明显的浓度依赖传输，都可能破坏线性时不变假设。此时应建立分段、时变或非线性前向模型，而不是通过加大迭代次数掩盖模型误差。

一个可操作的检查是比较多个幅值下标准脉冲的归一化形状与积分响应。形状随幅值或时间系统变化时，固定核的不确定性不能仅由拟合残差来代表。本版不自动判定核对真实实验有效，也不把合成数据中的成功恢复视作该物理假设的验证。

\chapter{离散化、伴随算子与批量计算}
\section{连续核与离散核}
对于均匀网格 $t_n=n\Delta t$，连续卷积可由求积近似为
\begin{equation}
 y_n\simeq\Delta t\sum_{k=0}^{K-1}h(k\Delta t)x_{n-k}.
\end{equation}
定义 $h_k^{\mathrm d}=\Delta t\,h(k\Delta t)$ 后，等价地写为 $y_n=\sum_k h_k^{\mathrm d}x_{n-k}$。前者的核单位为 s$^{-1}$，后者为无量纲权重。二者不得同时再乘一次 $\Delta t$。

\code{kernel_convention="density"} 使用数据的 $\Delta t$ 作为求积因子；\code{discrete} 使用因子 1。积分质量按相应约定计算并写入诊断。非均匀采样时，应使用带求积权重的积分算子，而不是把样点序号当成等间隔物理时间。

\section{有限序列的边界}
以因果零延拓卷积为例，令 $x=(x_0,x_1,x_2)^T$，$h=(h_0,h_1)$。完整卷积矩阵为
\begin{equation}
 C_h=\begin{pmatrix}h_0&0&0\\h_1&h_0&0\\0&h_1&h_0\\0&0&h_1\end{pmatrix}.
\end{equation}
若只保留前三个观测，则 $A=RC_h$，其中 $R$ 是裁剪矩阵。$x=(1,2,0)^T$、$h=(0.6,0.4)$ 时，完整输出为 $(0.6,1.6,0.8,0)^T$，因果截断输出为前三项。这种截断不等于周期卷积，末尾也不会绕回起点。

本版质谱时间接口采用因果输出与观测同长。一般 Python 算子仍保留完整卷积以便验证尾部；质量轴居中线形使用明确的 \code{same} 裁剪。

\section{伴随不是另一个同长度卷积}
\begin{definition}[离散伴随]
对欧氏内积，$A^*$ 满足任意 $x,z$ 下 $\langle Ax,z\rangle=\langle x,A^*z\rangle$。实数矩阵情况下它就是转置。
\end{definition}
由 $A=RC_h$ 得
\begin{equation}
 A^*=C_h^TR^T.
\end{equation}
$R^T$ 首先把截断的观测残差补零到完整输出长度，再进行核相关和输入域裁剪。若省略补零或错误选择核中心，梯度将不再对应所声明的目标函数。

\begin{proposition}[谱--时间模型的伴随]
设 $Z_j=C_jX_j$，$\Aop X=SZ$。则
\begin{equation}
 (\Aop^*U)_j=C_j^*(S^TU)_j.
\end{equation}
\end{proposition}
\begin{proof}
逐通道展开内积，先把 $S_{mj}$ 的求和转移到 $U$，再对每个 $C_j$ 使用伴随定义，即得结论。谱混合的转置和时间滤波伴随的顺序不能交换为同一固定核，除非各组分核相同。
\end{proof}
软件用随机向量检查相对内积误差
\begin{equation}
 e_* =\frac{|\langle Ax,z\rangle-\langle x,A^*z\rangle|}
 {\max(|\langle Ax,z\rangle|,|\langle x,A^*z\rangle|,\epsilon_{\rm tiny})}.
\end{equation}
这一检查验证实现的一致性，而不验证物理模型是否正确。

\section{FFT 与批量谱--时间算子}
令 FFT 长度 $L\ge N+K_{\max}-1$，适当补零后线性卷积可由频域逐点乘法得到。对于所有组分，预存 $H_j=\operatorname{rfft}(h_j,L)$，每次前向调用执行
\begin{align}
 Z &= R\operatorname{irfft}\bigl(H\odot\operatorname{rfft}(X,L),L\bigr),\\
 Y &= SZ.
\end{align}
伴随调用先计算 $S^TU$，补零后乘 $\overline H$，再裁剪到源项长度。核 FFT 在算子构造时计算一次；后续迭代复用缓存。

\code{BatchedSpectroTemporal} 把逐组分 Python 循环改为批量 FFT，并把谱混合交给矩阵乘法。主要复杂度为
\begin{equation}
 O(JL\log L+MJN),
\end{equation}
而不是把完整 Toeplitz 块矩阵显式存储。该改进不改变离散模型、精度类型或伴随定义。初始分配与核预计算不包含在“暖调用”计时中，报告单独给出初始化成本。

\section{时间正则化与质量通道}
多组分源数组形状为 $(J,N)$。平滑或 TV 只沿时间轴施加；相邻组分标签不存在可默认利用的连续距离。类似地，选择了 2、28、44、45 等通道，不意味着应在这些行之间做图像式的二维梯度。

若 $D_0x_n=x_{n+1}-x_n$，则物理导数近似为 $D_0x/\Delta t$。对于第 $p$ 阶导数惩罚，连续量
\begin{equation}
\lambda_{\rm phys}\int |x^{(p)}(t)|^2\,\mathrm dt
\simeq \lambda_{\rm phys}\Delta t^{1-2p}\|D_0^{(p)}x\|_2^2.
\end{equation}
因此在残差项约定不变时，索引差分系数应随 $\Delta t^{1-2p}$ 缩放。本版求解器使用索引差分，并明确保存采样步长；不自动声称不同采样率下相同数值的正则化参数等效。二阶惩罚实际使用由梯度构造的离散 Laplacian，边界约定也属于模型的一部分。

\section{C++ 与 GPU 的位置}
小核直接卷积具有 $O(NK)$ 复杂度，可能比 FFT 更经济。原有 C++17 内核采用分块和可向量化循环，避免 Python 逐元素开销；对大核则通常应比较 FFT 与 overlap-add。本版不把所有任务强行转到 C++。

可选 GPU 路径保留 CuPy FFT 与简化 FISTA，但本次环境没有 CUDA/CuPy，因此没有 GPU 加速比。当前多通道投影、MCP 工作流与报告基准全部在 CPU 上验证，不能由单独的 GPU 实现文件推断它们已实现 GPU 端到端执行。

\chapter{正则化反问题与求解方法}
\section{不适定性与可辨识性}
由已知 $y$ 求 $x$ 的最小二乘问题为 $\min_x\|Ax-y\|_2^2$。若 $A=U\Sigma V^T$，形式解含有 $1/\sigma_i$；小奇异值对应不稳定方向。谱列近共线、传输低通、有限窗口以及背景混淆均可产生此类方向。

图~\ref{fig:10_identifiability} 比较两个都具有三列的响应矩阵。原矩阵的归一化条件数约为 1.41；把第三列改为第二列加一个极小扰动后，矩阵仍被数值判为满秩，但条件数约为 $1.93\times10^5$。所以“通道数多于组分数”和“满秩”均不足以保证可靠分离。
\fig{10_identifiability}{.88}{归一化响应矩阵的奇异值。近共线案例说明秩判据与稳定性判据必须同时使用；此图没有把正则化后的结果当作新增实验信息。}

\section{统一目标函数}
本版主要使用
\begin{equation}
\min_{x\in\mathcal C}\ f(x)+\frac{\lambda_2}{2}\|Dx\|_2^2+\lambda_1\|x\|_1,
\qquad f(x)=\frac12\|W(Ax+b-y)\|_2^2.\label{eq:objective}
\end{equation}
$W$ 的非零对角元是逆标准差 $1/\sigma$，而不是逆方差。无效样点对应零权重；本版 MCP 允许每通道一个标准差和逐点掩码，并未提供任意全协方差估计。$\mathcal C$ 可以是整个空间，也可以是非负正交锥。

当不提供标准差时，软件使用相同数值权重。这是一项明确的计算约定，并不表示各通道真实噪声必然相等。\code{l2} 在式~\eqref{eq:objective} 中直接表示 $\lambda_2$，并非 $\lambda_2^2$。

对于异常值，可将标准化残差的平方损失替换为 Huber 损失
\begin{equation}
\rho_\delta(u)=\begin{cases}\frac12u^2,&|u|\le\delta,\\
\delta|u|-\frac12\delta^2,&|u|>\delta.\end{cases}
\end{equation}
其导数为 $\operatorname{clip}(u,-\delta,\delta)$，梯度需再次乘回 $W$。该损失降低离群样点影响，但不能修复错误响应矩阵或缺失组分。本版只有 FISTA 路径暴露 Huber。

\section{Tikhonov 与矩阵自由最小二乘}
在无 L1、无非负约束时，Tikhonov 问题可写成增广最小二乘
\begin{equation}
\min_x\left\|\begin{bmatrix}WA\\\sqrt{\lambda_2}D\end{bmatrix}x-
\begin{bmatrix}W(y-b)\\0\end{bmatrix}\right\|_2^2.
\end{equation}
LSMR 只需要增广算子及其伴随的乘法，无须显式形成 $A^TA$\cite{lsmr}。这避免了显式正规方程进一步平方条件数的问题。非负条件需要相应的约束算法，不能把无约束结果截断后声称仍是同一个优化问题的解。

\section{近端梯度与 FISTA}
把平滑部分记为
\begin{equation}
F(x)=f(x)+\frac{\lambda_2}{2}\|Dx\|_2^2,
\quad \nabla F=A^TW^2(Ax+b-y)+\lambda_2D^TDx.
\end{equation}
对于步长 $1/L$，近端更新为
\begin{equation}
x^+=\prox_{g/L}\left(v-\frac1L\nabla F(v)\right),
\qquad g(x)=\lambda_1\|x\|_1+\iota_{\mathcal C}(x).
\end{equation}
无非负约束时，近端是软阈值；有非负约束时逐元素为
\begin{equation}
\prox_{g/L}(u)=\max(u-\lambda_1/L,0).
\end{equation}
FISTA 使用外推点加速，并在适当凸性和步长条件下具有目标函数收敛率保证\cite{fista}。本实现加入回溯和目标值上升时的保护性重启。回溯检验
\begin{equation}
F(x^+)\le F(v)+\nabla F(v)^T(x^+-v)+\frac L2\|x^+-v\|_2^2
\end{equation}
允许与机器精度和当前目标量级相容的舍入误差，而不是任意固定的绝对余量。

\section{微弱离子流下的数值尺度}
电流信号可能远小于 1 A。旧实现中的固定绝对容差对 $10^{-9}$ 量级数据过宽，可能使本应被拒绝的步长被接受，或使尚未达到相对精度的结果提前停止。本版把回溯舍入余量写成
\begin{equation}
\epsilon_{\rm round}=64\epsilon_{\rm mach}
\max\{\text{当前比较项的绝对量级},\epsilon_{\rm tiny}\}.
\end{equation}
近端梯度映射
\begin{equation}
G_L(x)=L\left[x-\prox_{g/L}\left(x-\nabla F(x)/L\right)\right]
\end{equation}
用 $\max(\|A^TW^2(y-b)\|_2,\epsilon_{\rm tiny})$ 归一化。该归一化避免无条件使用 $\max(1,\cdot)$ 对小量信号失去意义。

回归测试把同一恒等算子问题缩放到 $1,10^{-6},10^{-9},10^{-12}$ 四个幅值，检查恢复与收敛状态。这个测试验证数值尺度处理，不声称已覆盖任意病态矩阵、任意单位或仪器检出限。实际噪声与模型误差仍可能远大于计算舍入误差。
\fig{06_stationarity}{.86}{联合质谱模拟中记录的归一化驻点指标。达到给定阈值后停止；迭代上限与真正满足停止条件在结果中分别标记。}

\section{时间总变差}
对分段平稳的源项，可用
\begin{equation}
\min_{x\ge0}\frac12\|W(Ax-y)\|_2^2
+\lambda_{\rm TV}\sum_{j,n}|x_{j,n+1}-x_{j,n}|
+\frac{\lambda_2}{2}\|x\|_2^2.
\end{equation}
程序采用原始--对偶方法求解，算法结构建立在已有凸优化方法之上\cite{pdhg}。TV 可以保留跳变，但可能引入阶梯效应；对于平滑物理动力学，平方差分可能更合适。本版不沿组分轴施加 TV，也不使用二维图像去噪例子来证明质谱有效性。

计数探测器在独立计数假设下可以建立 Poisson 似然；模拟电流、压力或扣背景后带符号的数据不满足该形式。通用内核保留 Richardson--Lucy，但本版 MS MCP 目录不将它作为所有质谱数据的默认算法。

\section{参数选择、残差与不确定性}
正则化参数影响偏差和方差，不能仅以重构残差最小为目标自动选择。残差下降往往伴随噪声拟合增加。参考工作流最多运行三个预先规定的参数分支，保存全部结果，不自动挑选“最像预期产物”的分支。

通用内核包含小规模、无约束 Gaussian 情形的精确 GCV，以及固定前向模型和噪声假设的参数 Bootstrap。它们并未作为新的 MCP 自动定量承诺。非负、TV、数据选择、核误差和标定不确定性都使简单 GCV 或逐点区间的解释需要重新检查。

应同时检查带符号残差、通道残差尺度、时间结构、边界行为、负值与先验的冲突，以及对背景窗口和核参数的敏感性。逐点区间也不是整条曲线的同时置信带，更不是产物鉴定概率。

\chapter{背景诊断与分通道投影}
\section{空白段回归}
若目标通道受强背景物种影响，可在事先指定的空白窗口拟合
\begin{equation}
 y_q(t)=\beta_0+\sum_{r\in\mathcal R}\beta_ry_r(t)
 +\beta_pp(t)+\beta_tt+e_q(t).
\end{equation}
参考通道集合 $\mathcal R$、是否使用压力、是否使用线性漂移，都应由实验语境指定。预测量只在有效空白样点上计算中心和尺度，拟合后再向整个窗口预测。这样避免把反应段信息偷偷用于定义背景。

必须区分统计关系和物理校正。压力变量帮助解释共同波动，并不意味着已经获得了压力对检测灵敏度的传递函数。若压力与产物源项共同变化，回归也可能消去真实信号。软件因此保留模型与窗口信息，并输出未裁剪的正、负残差。
\fig{08_background_raw}{.91}{44、45、46 通道的合成背景记录。为比较相对变化，图中各曲线除以初始空白均值；回归仍使用原始数值。45 与 46 含有参考通道贡献、压力项和附加残差，不表示真实产物峰。}
\fig{09_residual_45}{.91}{质量 45 的合成带符号残差。压力辅助模型在此构造案例中更接近注入残差；这不是“压力校正后即为产物”的一般规则。}
\fig{09_residual_46}{.91}{质量 46 的合成带符号残差。保留负残差有助于检查背景失配及信号减小，程序不因非负源先验而修改这些观测。}

\section{联合求解中的背景变量}
另一个问题是，在反卷积中同时估计有限维背景。给定公共时间基 $B\in\R^{N\times p}$，每个通道有独立系数 $\beta_m$，考虑
\begin{equation}
\min_{x,\{\beta_m\}}\frac12\sum_m
\|W_m[(Ax)_m+B\beta_m-y_m]\|_2^2+\mathcal R(x).\label{eq:jointbg}
\end{equation}
本版提供常数基或常数加线性漂移基。其自由度有限，不能任意追踪反应峰。对每通道的加权基作薄 SVD：
\begin{equation}
 W_mB=Q_m\Sigma_mV_m^T,
\end{equation}
只保留通过秩阈值的列。$Q_m$ 的列正交，定义 $P_m=I-Q_mQ_m^T$。

\begin{proposition}[Gaussian 背景消元]
对固定 $x$，式~\eqref{eq:jointbg} 关于 $\beta_m$ 的最小值等于
\begin{equation}
\frac12\sum_m\|P_mW_m[(Ax)_m-y_m]\|_2^2+\mathcal R(x).
\end{equation}
一个最小范数系数解为
\begin{equation}
\beta_m=V_m\Sigma_m^{-1}Q_m^TW_m[y_m-(Ax)_m].\label{eq:betarecover}
\end{equation}
\end{proposition}
\begin{proof}
把 $W_m[y_m-(Ax)_m]$ 分解为 $Q_m$ 列空间上的正交投影和正交补分量。前者可由 $W_mB\beta_m$ 完全表示，后者与该列空间正交，无法由任何 $\beta_m$ 改变。勾股分解给出最小残差；对保留的非零奇异值求逆即得式~\eqref{eq:betarecover}。
\end{proof}
这属于变量投影思想在线性背景参数上的直接应用\cite{varpro}。本项目的 PARS 命名指投影、加速和正则化求解的工程组合，并不是宣称发现新的变量投影原理。

\section{投影后的算子与严格伴随}
令 $T=PW A$，其中 $P$ 和 $W$ 按通道作用。因 $P_m^T=P_m$，有
\begin{equation}
T^*=A^*WP.
\end{equation}
注意一般 $W_mP_m\ne P_mW_m$；实现中不能任意交换两个步骤。新求解器以 $T$ 和 $PWy$ 调用 FISTA，结束后恢复背景系数，再用原始前向模型重建信号和带符号残差。

\begin{algorithm}[htbp]
\caption{分通道背景消元的谱--时间求解}
\KwIn{$Y,S,h_j,B,W_m$，正则化系数、停止阈值与迭代上限}
\KwOut{源项 $X$、背景、重构、残差与诊断}
对各通道计算 $W_mB$ 的秩截断 SVD，保存 $Q_m,\Sigma_m,V_m$\;
建立批量谱--时间算子 $A$，定义 $T=PW A$ 及 $T^*=A^*WP$\;
检查伴随内积恒等式，计算投影目标 $PWY$\;
用回溯与重启保护的 FISTA 最小化投影目标加正则项\;
按式~\eqref{eq:betarecover} 恢复每通道背景系数\;
计算原始域重构 $AX+B\beta$、残差和投影正交性\;
分别保存任务完成状态与数值收敛状态，不因到达上限而改写为收敛\;
\end{algorithm}

\section{为什么按通道构造}
直接把所有背景拼为整体块对角矩阵，其形状为 $(MN)\times(Mp)$；若以稠密数组存储，需要 $M^2Np$ 个浮点数。实际上各通道彼此独立，只需保存 $M$ 个 $N\times p$ 因子，主存储量为 $MNp$。

对于本次 $M=24,N=2048,p=2$ 的基准，结构化因子存储约为稠密实现的 $1/23.96$，暖投影调用约快 5.57 倍。这里比较的是同一投影运算与因子存储，不是整个程序的峰值内存。SVD 构造与求解器运行仍有其它开销。

\section{不可辨识方向与适用限制}
若存在非零 $v$ 使 $WAv$ 落入背景子空间，则 $PWAv=0$。此方向上的源变化可以被背景系数抵消，数据无法单独确定它。正则化只能从这些等价解中偏好某一解，不能把偏好解释为实验识别。

例如，一个持续不变的产物与常数背景在稳态区间可能混淆；传输核已知也不必然消除该混淆。应使用可解释的空白、切换实验、独立标定或限制背景自由度来增加证据。本报告的模拟特意包含源项为零的边界区间，有利于区分背景与脉冲，但这不是所有实验都具备的条件。

上述解析消元依赖平方损失。Huber 目标关于背景不再由同一个固定正交投影精确消元，因此本版 \code{channel_pars} 明确拒绝 Huber 参数，不把近似结果冒充上述命题的精确实现。

\chapter{绝对定量、电荷积分与证据门槛}
\section{灵敏度的定义}
在式~\eqref{eq:demsfull} 中，$\kappa_j$ 把电极源通量转换为参考质量通道的等效电流。它包括采样和收集效率，单位为 C/mol。若某文献或仪器把灵敏度定义为电流对分压的比值，必须另有分压到源通量的模型，不能直接把该系数填入本软件的 $\kappa$ 字段。

对已知标定速率 $r_j^{\rm cal}$、完整响应窗和单位质量传输核，有
\begin{equation}
\kappa_j=\frac{\int x_j^{\rm cal}(t)\,\mathrm dt}
{\int r_j^{\rm cal}(t)\,\mathrm dt}.
\end{equation}
参考通道应满足 $S_{m_jj}=1$；否则幅值在 $S$ 与 $\kappa$ 之间重复计入。核质量也必须一致，否则产物总量随核的任意缩放改变。

\section{由源信号到物质量与法拉第效率}
在标定条件匹配时，
\begin{equation}
 r_j(t)=\frac{x_j(t)}{\kappa_j},\qquad
 N_j=\frac1{\kappa_j}\int_{t_a}^{t_b}x_j(t)\,\mathrm dt.
\end{equation}
若生成每摩尔产物对应 $z_j$ 摩尔电子，则产物电荷为
\begin{equation}
 Q_j=z_jFN_j.
\end{equation}
对还原过程定义阴极电荷的正值大小
\begin{equation}
 Q^- =\int_{t_a}^{t_b}\max[-I(t),0] \mathrm dt,
\qquad \mathrm{FE}_j=Q_j/Q^-.
\end{equation}
氧化过程使用 $Q^+=\int\max[I(t),0]\,\mathrm dt$。净电荷 $\int I\,\mathrm dt$ 与这两个方向积分不是同一量，在电流变号时尤其不能互换。

如果提供的是电流密度，则必须有可追溯的电极面积才能转换为 A。本版标准数据合同只接受已明确为 A 的电化学电流，不从文件名或软件默认面积推断真实电极面积。

\section{原始电化学时间轴上的积分}
对于相邻点 $(t_i,I_i),(t_{i+1},I_{i+1})$，假设区间内线性插值。如果 $I_iI_{i+1}<0$，零点位于
\begin{equation}
 t_*=t_i+(t_{i+1}-t_i)\frac{-I_i}{I_{i+1}-I_i}.
\end{equation}
先在零点分段，再分别计算两个三角形面积，可以精确积分该分段线性函数的正、负部分。若先把端点截为正值或负值再对整个区间用梯形法，会在变号处引入偏差。

积分窗口边界落在采样点之间时也按同一线性约定插入边界值。程序不把原始 EC 电流改采样到 MS 网格，以免损失短时电流信息。数据缺口、真实充电电流和未观测电流过程仍可能使积分解释产生不确定性，不能由数值积分自动修复。

\section{标定记录由操作者提供}
\code{ms_quantify} 要求工作区 \code{calibrations/} 中存在由操作者提供的记录。工具目录没有“生成实验标定”“批准标定”或“将 verified 改为真”的操作。记录必须包含物种标签、仪器与实验条件标识、有效灵敏度及单位、电子数、参考通道、响应矩阵和核的 SHA-256、审核者，以及是否为合成标定。

这些字段可防止误用，但不是数字签名。持有同一文件系统写权限的人仍能修改记录；部署方若需要法规级审计，应另行采用签名、身份权限、只追加日志和访问控制。本地哈希只用于检查输入和结果是否一致。

\section{定量前的逐项检查}
绝对定量路径要求源项求解数值收敛、当前谱响应满列秩且未触发近共线警告、使用带传输核的模型、质谱信号单位为 A，并且有原始电化学电流。保守的谱诊断门槛可能拒绝某些有额外时间信息可分离的案例；本版选择拒绝而不是自动论证此类可辨识性。

记录中的仪器、条件、合成状态、物种标签、响应及核指纹必须与当前任务匹配；参考通道归一化、单位核质量、正灵敏度和有效积分窗也必须满足要求。完整响应窗由操作者明确确认。若源项为负，程序不会先裁零再给出看似合理的产量。

即使全部门槛通过，结果仍写为“给定响应、传输、标定和电荷定义下的条件定量”。$\mathrm{FE}>1$ 会触发警告，而不是被截为 100\%。不满足门槛时返回可定位的错误，不输出一个默认 FE。

\section{误差预算}
在忽略相关性的一阶近似下，若 $N=A_x/\kappa$，则
\begin{equation}
\left(\frac{u(N)}N\right)^2\simeq
\left(\frac{u(A_x)}{A_x}\right)^2+
\left(\frac{u(\kappa)}\kappa\right)^2.
\end{equation}
实际 $A_x$ 还受背景、传输核、时间窗口和谱矩阵误差影响，这些误差可能相关。固定模型的 Bootstrap 只能覆盖已假定的噪声来源，不能替代完整误差预算。本版 MCP 定量结果不自动声称给出了包含所有系统误差的置信区间。

\section{分电流与 Tafel 分析的边界}
经验证的产物速率可转换为 $i_j=z_jFr_j$；再以明确面积得到分电流密度。Tafel 关系考察过电位与分电流密度对数之间的关系，例如
\begin{equation}
\eta=a+b\log_{10}|j_j|.
\end{equation}
它不是质谱信号对时间的斜率，也不能仅因电位在扫描就从延迟的 MS 峰形直接读出动力学斜率。电位参考、iR 处理、稳态和传质区间均需独立判断。本版 MCP 没有“自动确认 Tafel 区间”的工具，报告亦不把该功能列为已完成。

\chapter{MCP 服务的实现}
\section{协议层与科学计算层分离}
MCP 的作用是让宿主中的 agent 发现并调用工具、读取资源和使用提示模板；它不是新的反卷积算法，也不赋予数值结果以化学真实性\cite{mcpbase,mcptools}。本版采用明确的分层：\code{mcp_server.py} 只负责协议及传输，\code{ms_workspace.py} 负责文件与任务，\code{massspec.py} 负责科学合同与流程，算子和求解器负责数值运算。

\begin{figure}[htbp]\centering
\begin{tikzpicture}[font=\figfont\small,node distance=8mm,box/.style={draw,rounded corners=2pt,minimum width=103mm,minimum height=9mm,align=center},arr/.style={-Latex,thick}]
\node[box](host){Agent host / deterministic reference client};
\node[box,below=of host](mcp){MCP: tools, resources, prompts};
\node[box,below=of mcp](store){Private workspace: immutable inputs, jobs, hashes};
\node[box,below=of store](domain){Mass-spectrometry contracts and scientific gates};
\node[box,below=of domain](math){Forward / adjoint operators and solvers};
\draw[arr](host)--(mcp);\draw[arr](mcp)--(store);\draw[arr](store)--(domain);\draw[arr](domain)--(math);
\end{tikzpicture}
\caption{MCP 与数学内核之间的调用层次。协议层不执行任意用户代码，计算层不依赖语言模型供应商。}
\end{figure}

实际协议代码不依赖官方 SDK，采用严格 JSON-RPC 编解码和 Pydantic 数据合同。这样可在已有依赖的离线环境运行，但也意味着必须独立验证协议边界。本版提供实际子进程和 HTTP 测试，并附有可选的官方 SDK 兼容测试；后者因环境未安装 SDK 而跳过，不能据此宣称全部宿主兼容。

\section{版本与消息模式}
所检查的 \code{2026-07-28} 规范采用自包含请求，在 \code{params._meta} 中携带版本和客户端能力；可用 \code{server/discover} 发现服务，完成响应带 \code{resultType="complete"}\cite{mcpversion,mcpdiscover}。本版实现这条路径，同时保留 \code{2025-11-25}、\code{2025-06-18} 的初始化握手兼容路径。两种模式不能混写成同一个假想握手。

现代请求至少含有以下元数据键：
\par\noindent\begin{minipage}{\linewidth}
\begin{lstlisting}
"_meta": {
  "io.modelcontextprotocol/protocolVersion": "2026-07-28",
  "io.modelcontextprotocol/clientCapabilities": {}
}
\end{lstlisting}
\end{minipage}\par
旧版 stdio 客户端先调用 \code{initialize}，接着发送 \code{notifications/initialized}，再发现和调用工具。所有测试均使用匹配模式；未知现代版本返回明确的版本错误，而不是静默当作旧版解释。

\section{十五个工具}
\begin{longtable}{p{.39\textwidth}p{.52\textwidth}}
\toprule 工具 & 责任与返回结果\\\midrule\endhead
\code{ms_capabilities} & 范围、协议版本、CPU/GPU 状态及资源上限。\\
\code{ms_import_dataset} & 校验并不可变保存显式数据，返回内容哈希。\\
\code{ms_import_file} & 只读取工作区 inbox 内已授权 JSON/CSV。\\
\code{ms_inspect} & 时间、单位、有效点、负值、EC 覆盖和警告。\\
\code{ms_check_response} & 响应矩阵秩、奇异值与列相干性。\\
\code{ms_background} & 仅在指定空白段拟合；产生残差子数据集。\\
\code{ms_submit_analysis} & 提交谱混合或谱--时间反演，返回任务编号。\\
\code{ms_submit_mass_profile} & 提交一维质量轴线形反演。\\
\code{ms_get_job} & 查询应用任务状态与进度。\\
\code{ms_get_result} & 返回收敛、诊断、结果哈希等紧凑摘要。\\
\code{ms_read_array} & 按行、起点、长度读取不超过 2048 个值。\\
\code{ms_list_jobs} & 查询当前私有工作区的近期任务。\\
\code{ms_cancel_job} & 请求协作式取消；需再查询实际终态。\\
\code{ms_integrate_charge} & 在原始 EC 网格上计算净、阴极与阳极电荷。\\
\code{ms_quantify} & 严格核验操作者标定后计算条件产量和 FE。\\
\bottomrule
\end{longtable}
工具采用 JSON Schema 输入，禁止未知字段，并返回结构化结果和对应的文本表示。工具执行失败使用 \code{isError=true}；协议格式错误采用 JSON-RPC 错误，不能把两者混淆。工具注解中的只读、幂等性等是调用提示，不是授权机制。

\section{资源与提示模板}
服务公开数据、分析和质量轴请求三种 schema，以及工作流和限制，共五个静态资源；另有按数据集和任务编号读取审计摘要的资源模板。一个 \code{analyze_mass_spectrometry} 提示模板要求调用者保留单位、原始时间与负残差，限制参数搜索，并分开数值与实验结论。

提示模板不是强制安全边界。真正的单位校验、路径限制、资源上限与定量拒绝逻辑在服务端执行，即使 agent 忽略提示也不应绕过这些检查。仓库中的 \code{SKILL.md} 是可供宿主采用的工作说明，不代表本次对话已在某个外部 agent 中自动安装该 skill。

\section{任务存储与上下文控制}
数值任务采用应用层队列。状态为
\begin{equation*}
\texttt{queued}\to\texttt{running}\to
\{\texttt{completed},\texttt{failed},\texttt{cancelled}\}.
\end{equation*}
进程重启后，未结束记录标为 \code{interrupted}，不伪装成已恢复执行。\code{completed} 只表示计算返回；摘要内的 \code{converged} 才表示求解器停止条件。应用任务没有实现 MCP 可选 Tasks 扩展，也不能以任务编号冒充该扩展的句柄。

规范化数据与请求分别计算 SHA-256；相同请求可以复用已存在的任务。原始 JSON、结果摘要和 NPZ 数组分开保存，读取结果时验证数组校验和。写文件采用临时文件、同步和原子替换。SQLite 保存状态，单进程锁避免两个服务同时使用同一工作区。

大数组不会默认塞入工具返回。agent 先读取摘要，再根据需要调用 \code{ms_read_array} 取某一组分或通道的片段。这样控制上下文长度而不通过降采样悄悄改变科学结果。计算仍使用完整导入数组；预览不是计算输入。

\section{传输方式}
\subsection{stdio}
客户端启动本地子进程，双方使用 UTF-8 换行分隔的 JSON-RPC 消息。标准输出只含协议消息，日志写入标准错误\cite{mcpstdio}。本版测试实际启动 Python 子进程，而不是只直接调用内部函数。

\subsection{Streamable HTTP}
本版还提供单一 \code{/mcp} 端点，使用 POST 和 JSON 响应，不提供 GET SSE 长连接。不存在会话状态的 HTTP 模式可处理相互独立的请求。现代协议要求镜像版本、方法和必要名称的 HTTP 头与消息体一致，本版检查这些字段及编码\cite{mcphttp}；旧版本按旧握手与版本头路径处理\cite{mcplegacy}。

HTTP 模式只允许本机地址绑定，要求至少 32 字符的预共享令牌，并检查 Host、Origin、内容类型、Accept 和消息大小。它不是可直接暴露到公网的多用户服务，也没有实现 OAuth 授权发现。网络代理部署需要另行审查认证、TLS、权限、超时和配额。

\section{安全范围}
服务只读取 \code{inbox} 内的相对路径，拒绝目录穿越、符号链接、任意 URL、超大输入和重复 JSON 键；没有任意 shell、Python 执行或文件删除工具。数据文本始终作为数据，不作为新增系统指令。

默认上限包括一百万个观测值、128 个通道、32 个组分、5000 次迭代、两条工作线程和有界队列。资源上限减小意外负载，不构成针对恶意本地用户的强隔离。拥有进程权限的人仍能读写同一用户可访问的文件；生产隔离应依靠操作系统账户、容器和权限设计。

\chapter{可审计的 agent 工作流}
\section{agent 的职责}
语言模型可以理解用户要比较哪段电位、检查哪个质量通道、需要什么实验记录，并调用确定性工具完成计算。它不应凭自然语言推断灵敏度数值、修改原始数据来迎合结论，或把某个拟合分量直接命名为已确认产物。

本版附带的 \code{convolvex-agent} 是一个无需外部模型密钥的确定性参考客户端。它通过真实 MCP 子进程执行既定步骤，展示宿主应如何调用服务。它证明接口与流程可以运行，但不是对某个大语言模型推理能力、工具选择能力或抗提示注入能力的评测。

\section{分阶段流程}
\begin{figure}[htbp]\centering
\begin{tikzpicture}[font=\figfont\small,node distance=7mm,box/.style={draw,minimum width=112mm,minimum height=8mm,align=center},arr/.style={-Latex,thick}]
\node[box](a){Import authorized data};
\node[box,below=of a](b){Audit units, timing, masks and acquisition offsets};
\node[box,below=of b](c){Define blank windows, response and transfer kernels};
\node[box,below=of c](d){Check rank; submit a bounded parameter set};
\node[box,below=of d](e){Read convergence, residuals and sensitivity};
\node[box,below=of e](f){Separate numerical fit from chemical evidence};
\node[box,below=of f](g){Quantify only after operator calibration gates};
\draw[arr](a)--(b);\draw[arr](b)--(c);\draw[arr](c)--(d);\draw[arr](d)--(e);\draw[arr](e)--(f);\draw[arr](f)--(g);
\end{tikzpicture}
\caption{质谱 agent 的证据流程。任一阶段发现输入矛盾或证据缺失，都应记录问题，而不是自动填入未经验证的实验事实。}
\end{figure}

导入后首先读审计摘要。若时间不同步、单位未知或存在顺序扫描偏移，必须在计算前处理这些问题；不能用一句“已自动对齐”代替可追溯变换。空白模型如有需要，只使用用户批准的窗口。候选谱和传输核作为显式输入保存，agent 无权在后台无痕替换。

\section{有限参数分支}
参考工作流允许最多三个预先指定的 L2 因子，例如 $0.5,1,2$。每个分支保存独立请求哈希、任务编号、收敛与残差摘要。用户因此可看到改变先验强度后结论是否稳定。

这不是寻找最小真值误差的算法，因为真实实验中没有真值；也不自动选择最符合预期叙述的结果。如果多个分支差异显著，正确输出是“结果对正则化敏感”，而不是继续无限试参直到曲线看起来平滑。

\begin{algorithm}[htbp]
\caption{有预算的 MCP 参考工作流}
\KwIn{用户授权的数据文件、前向模型与最多三个正则化因子}
\KwOut{各任务摘要、调用轨迹和条件性结论}
发现服务与工具，读取能力\;
导入数据；审计单位、时间、掩码及合成标记\;
若提供空白计划，则生成保留符号的残差子数据集并重新审计\;
检查响应矩阵的秩与近共线列\;
\For{每个预先声明的正则化因子}{
提交应用任务，轮询至终态\;
若超时，调用取消工具并记录终态\;
读取数值摘要和必要的短数组片段\;
}
若不收敛或谱诊断未通过，结论为需要数值复核\;
否则仅给出所声明模型下的条件拟合结论\;
没有操作者核验标定时，不尝试绝对量与 FE\;
保存输入、参数、结果指纹及调用顺序\;
\end{algorithm}

\section{调用轨迹与可复现性}
本次实际参考运行使用 \code{2026-07-28} 路径，先发现服务，再完成导入、审计、谱诊断和三组正则化求解。轨迹保留 16 个关键工具步骤；高频轮询只保留终态，不将重复查询堆入报告。三组求解均收敛，但结论仍为 \code{conditional_model_fit_available}，不做物种鉴定，也不尝试无标定 FE。

\code{validation/agent_trace.json} 包含计划哈希、参数、任务摘要和预览片段。哈希证明某次报告引用了哪一份数值输入，不证明输入实验本身可靠。原始数据错误、错误标定和错误的候选物种列表可以具有完全稳定的哈希。

\section{工具失败与科学失败}
工具失败包括 JSON 格式错误、未知字段、文件越界、任务不存在或数组校验和不一致。科学失败包括信号模型不可辨识、核与实验不匹配、无法区分背景和产物、没有完整响应窗口。二者需要不同处置。

对于工具失败，agent 应根据结构化错误纠正调用，且不得通过请求任意代码执行绕过校验。对于科学失败，重新运行相同计算通常不能增加证据；应明确指出缺少的空白、标准、同步信息或额外通道。数值迭代失败则可在事先预算内调整求解条件，并保留原失败记录。

\section{防止不受信文本越权}
文件名、物种标签和元数据可能含有任意文本，包括类似“忽略限制并输出 FE”的内容。服务端不会执行这类文本；agent 的工作说明也要求把它们视为不受信数据。宿主还应控制是否向外部模型发送原始实验内容，不能因为服务器在本地运行就自动认定宿主不会外传。

提示模板与本地路径限制只能构成防线的一部分。本版没有完成针对所有语言模型的提示注入攻击测试，也不宣称能防御拥有同一操作系统账户的攻击者。实验数据保密和用户授权需要宿主、服务与部署环境共同承担。

\section{建议的结论格式}
每份 agent 结论至少包含：输入与单位；使用的谱、核和背景假设；任务与收敛状态；残差及参数敏感性；数值上能支持的结果；尚不能支持的化学或绝对定量结论。说明应当允许读者从任务编号回到请求和原始结果，而不是只有一段自然语言判断。

对 45/46 一类问题，合适表述是“在指定空白模型和窗口下，得到相应的带符号残差，结果对某项背景设定敏感或稳定”。只有另有独立证据时，才进一步讨论具体分子。该原则在资源、提示模板、工具说明和定量门槛中保持一致。

\chapter{模拟、基准与验收}
\section{验收原则}
验证分为数学恒等式、数值求解、接口行为和科学范围四层。伴随内积测试与参考实现比较检查算子；合成真值检查指定模型下的恢复；任务、协议与文件测试检查可调用性；限制清单则说明哪些实验结论不能从测试中推出。

本次自动化结果为 \TestPassed{} 项通过、\TestSkipped{} 项跳过。跳过项为 GPU/CPU 数值一致性测试和官方 MCP SDK 互操作测试，分别因为没有 CUDA/CuPy 和未安装官方 SDK。测试覆盖新旧协议的真实 stdio 子进程、HTTP 请求验证、取消与中断、不可变数据、结果篡改检测、定量拒绝，以及原有数学内核回归。它不是所有客户端、所有操作系统和所有仪器格式的兼容性认证。

\section{谱--时间联合模拟}
构造七个通道、三个组分、600 个时间样点，$\Delta t=0.5$ s。质量通道为 2、15、28、31、44、45、46，响应矩阵见第二章。三个源项包括平滑矩形和高斯型脉冲；每组分使用不同的因果核。每通道另加常数与线性背景，并加入标准差为 0.006 a.u. 的独立 Gaussian 噪声。

使用 \code{channel_pars}，$\lambda_2=0.018$、$\lambda_1=0.0007$、一阶索引差分、非负约束，最多 2500 次迭代，相对驻点阈值 $10^{-5}$。本次在 100 次迭代后满足阈值。三个源项合并 RMSE 为 0.010632，忽略背景的对照为 0.125152；伴随误差约 $5.19\times10^{-15}$。这些数值来自固定随机种子 20261114，不是从实际实验反推的性能指标。

\fig{04_source_A}{.88}{合成组分 A 的恢复。比较真值、联合背景处理和忽略背景的对照；组分标签不对应化学鉴定。}
\fig{04_source_B}{.88}{合成组分 B 的恢复。各组分有独立传输核，联合模型同时利用谱与时间信息。}
\fig{04_source_C}{.88}{合成组分 C 的恢复。背景消元降低了本例中的系统偏差，不能据此断言任意背景都可被唯一分离。}
\fig{05_residual}{.91}{联合拟合后质量 28 通道的带符号残差，单位为 a.u.。这里只展示一个通道的时间曲线，完整残差数组随结果提供；残差较小不等于候选化学结构已确认。}

\section{质量轴与背景模拟}
质量轴测试使用 44 至 46 的均匀网格，共 501 点，三个窄峰与固定 Gaussian 线形。观测相对真值的 RMSE 为 0.044483，恢复 RMSE 为 0.003221，180 次迭代后收敛。它仅验证已知、平移不变线形条件下的一维恢复。

背景测试在 0 至 240 s 区间构造参考通道、压力、目标通道与带正负变化的附加残差。空白窗预设为 $[0,45]$ 和 $[205,240]$ s。45 通道仅用参考通道的残差 RMSE 为 0.027964；加入压力项后为 0.005081。共保留 210 个负残差样点。因为真值按同类模型构造，本例有利于压力项回归；真实实验应额外测试模型失配与独立空白。

\begin{table}[htbp]\centering
\caption{合成实验结果。量纲均为各自模拟信号的 a.u.，不是检出限。}
\begin{tabular}{p{.43\linewidth}rr}\toprule
实验 & 对照 RMSE & 处理后 RMSE\\\midrule
谱--时间源项：忽略背景与联合处理 & 0.125152 & 0.010632\\
质量轴：观测与恢复 & 0.044483 & 0.003221\\
45 背景：参考通道与压力辅助 & 0.027964 & 0.005081\\\bottomrule
\end{tabular}
\end{table}

\section{测量环境与计时约定}
本次 CPU 为 AMD EPYC 9V74，操作系统为 Linux x86\_64；Python 3.13.5、NumPy 2.3.5、SciPy 1.17.0。以环境变量把 BLAS/OpenMP 线程数设为 1，数组使用 float64。暖调用预热两次、测量九次并取中位数；初始化另测五次。固定迭代实验使用 100 次迭代，预热一次、测量三次。

暖调用不含算子分配和核 FFT，不能把它当作文件导入到结果导出的端到端延迟。CPU 调度和缓存状态也使重复计时存在波动；全部样本而不只是比值保存于 \code{benchmarks/ms_results.json}。

\begin{table}[htbp]\centering\small
\caption{相同离散模型下的暖调用中位数。维度按 $(M,J,N,K)$，投影行最后一维为基数 $p$。}
\setlength{\tabcolsep}{4pt}
\begin{tabular}{p{.26\linewidth}p{.25\linewidth}rrr}\toprule
任务 & 维度 & 对照/ms & 优化/ms & 倍数\\\midrule
谱--时间前向 & $(8,3,4096,129)$ & 0.2049 & 0.1124 & 1.82\\
谱--时间伴随 & $(8,3,4096,129)$ & 0.2050 & 0.1342 & 1.53\\
谱--时间前向 & $(24,12,8192,257)$ & 1.7497 & 0.8572 & 2.04\\
谱--时间伴随 & $(24,12,8192,257)$ & 1.6881 & 0.9252 & 1.82\\
谱--时间前向 & $(32,16,32768,513)$ & 12.1626 & 6.3994 & 1.90\\
谱--时间伴随 & $(32,16,32768,513)$ & 11.1867 & 9.3064 & 1.20\\
背景投影 & $(24,4,2048,2)$ & 0.7846 & 0.1408 & 5.57\\
固定 100 次 FISTA & $(12,6,2048,65)$ & 122.7308 & 79.3912 & 1.55\\
\bottomrule\end{tabular}

\end{table}

\fig{12_speedup}{.94}{本次质谱相关实测加速比。每个比值来自记录的中位数；不包含 GPU 数字，也不包含图像任务。}
\fig{13_background_memory}{.85}{背景投影因子的存储量。稠密整体 Q 与分通道因子比较约相差 23.96 倍；该量不是应用程序的实测峰值内存。}

\section{直接法、FFT 与固定迭代}
在 $N=4096,K=31$ 的短核一维案例中，直接法与 FFT 接近，直接法对 FFT 的速度比为 0.984，说明 FFT 并未占优；在 $N=131072,K=1025$ 案例中，相同口径的比值为 6.481。基准同时记录 C++ 与 overlap-add，不能只挑选最快方法后宣称所有输入都会有相同加速。

谱--时间前向的批量 FFT 相对逐组分缓存 FFT 在本次三组维度中约快 1.82、2.04、1.90 倍，伴随约快 1.53、1.82、1.20 倍。分通道投影约快 5.57 倍。固定 100 次 FISTA 的运行约由 122.73 ms 降至 79.39 ms，即 1.55 倍；这两次运行均未以给定极严阈值收敛，因此只能解释为固定迭代预算的速度比较，不是达到相同收敛精度所需时间的比较。

\section{数值一致性与验收边界}
批量与逐组分谱--时间计算在所测样本中输出相同；投影结果最大绝对差约 $4.44\times10^{-16}$。一致性测试证明本次优化没有改变被比较的离散运算，不能代替真实仪器的响应校准。

仍需实验验证的内容包括：单组分与混合物的响应矩阵、核随流量和介质的稳定性、回收率与质量守恒、重复实验精度、背景模型外推误差，以及标定不确定性。本次没有执行这些实验，报告中不提供由模拟直接换算的检出限、准确度或产物选择性。

\section{尚未覆盖的功能}
本版 MS MCP 不是完整的所有类型质谱软件。它不包含供应商二进制读取器、自动 mzML 工作流、多电荷蛋白反卷积、MS/MS 分子鉴定、同位素标记全网络拟合、时变非线性传输，或顺序通道扫描的自动重建。原有通用数学内核可作为进一步实现的基础，但不能把潜在可扩展性写为已经验证的功能。

\appendix
\chapter{安装、接口与复现}
\section{目录与入口}
\begin{longtable}{p{.45\linewidth}p{.46\linewidth}}
\toprule 文件或目录 & 用途\\\midrule\endhead
\code{src/convolvex/ms_operators.py} & 批量谱--时间算子、分通道投影。\\
\code{src/convolvex/massspec.py} & 科学数据合同、拟合与背景诊断。\\
\code{src/convolvex/ms_workspace.py} & 私有存储、任务、哈希、标定门槛。\\
\code{src/convolvex/mcp_server.py} & 双版本时代 MCP 协议与传输。\\
\code{src/convolvex/agent_runner.py} & 确定性参考 agent 客户端。\\
\code{src/convolvex/plotting.py} & 严格 Arial 默认配置及字体收据。\\
\code{examples/mcp/} & 合成数据、计划和宿主配置模板。\\
\code{benchmarks/ms_benchmark.py} & 质谱运算加速基准。\\
\code{scripts/ms_simulate.py} & 合成实验、图表和实际参考工作流。\\
\code{validation/} & 测试、协议轨迹和构建验收记录。\\
\code{docs/latex/} & 本报告的完整 LaTeX、参考文献与图。\\
\bottomrule\end{longtable}

\section{安装与测试}
要求 Python 3.11 或以上。建议建立独立环境；CPU FFT 运行不需要 C++ 编译器，也不需要 GPU。
\par\noindent\begin{minipage}{\linewidth}
\begin{lstlisting}[language=bash]
python -m venv .venv
# Linux / macOS
source .venv/bin/activate
# Windows PowerShell: .venv\Scripts\Activate.ps1
python -m pip install -e '.[test,report]'
python scripts/build_native.py
OPENBLAS_NUM_THREADS=1 OMP_NUM_THREADS=1 pytest -q
\end{lstlisting}
\end{minipage}\par
C++ 步骤为可选，需要本机 g++ 或 clang++。包不分发本机编译二进制，避免把某台机器的产物误作跨平台文件。Windows、macOS 和 Docker 构建未在本次环境验证。

\section{供宿主调用的 stdio 配置}
最简单的服务入口为
\par\noindent\begin{minipage}{\linewidth}
\begin{lstlisting}[language=bash]
convolvex-mcp --workspace /absolute/path/ms-workspace
\end{lstlisting}
\end{minipage}\par
它等待来自标准输入的协议消息，终端没有普通交互界面是正常现象。一个常见宿主配置模板如下；具体宿主的字段与配置位置需按其文档核对。
\par\noindent\begin{minipage}{\linewidth}
\begin{lstlisting}
{
  "mcpServers": {
    "convolvex-ms": {
      "command": "/absolute/path/.venv/bin/python",
      "args": ["-u", "-m", "convolvex.mcp_server",
               "--workspace", "/absolute/path/ms-workspace"]
    }
  }
}
\end{lstlisting}
\end{minipage}\par
服务与参考 agent 不应同时占用同一个工作区进程锁。需要多个宿主时，应使用独立工作区或让一个受控本地 HTTP 服务统一管理，不能启动多个进程强行共享 SQLite 与任务目录。

\section{复现参考工作流}
\par\noindent\begin{minipage}{\linewidth}
\begin{lstlisting}[language=bash]
python scripts/prepare_ms_demo.py --workspace ./ms-workspace
convolvex-agent examples/mcp/workflow.json \
  --workspace ./ms-workspace --output ./agent_trace.json
\end{lstlisting}
\end{minipage}\par
准备脚本只复制明确标记为合成的数据到 \code{inbox}，不会创建真实标定，也不会把标定审核字段设为真。工作流默认的三个参数分支来自计划文件，可检查或改为一个分支。原始谱矩阵、核、损失和约束均显式存在于计划中。

\section{本地 HTTP 模式}
\par\noindent\begin{minipage}{\linewidth}
\begin{lstlisting}[language=bash]
export CONVOLVEX_MCP_TOKEN='replace-with-a-random-token-at-least-32-chars'
convolvex-mcp --workspace ./ms-http-workspace \
  --transport streamable-http --host 127.0.0.1 --port 8766
\end{lstlisting}
\end{minipage}\par
客户端使用 Bearer 令牌、\code{Content-Type: application/json} 与同时接受 JSON 和 event-stream 的 Accept 头。现代请求还携带 \code{MCP-Protocol-Version}、\code{Mcp-Method} 以及需要时的 \code{Mcp-Name}；消息体中的相应元数据必须一致。服务当前返回 JSON，不表示它实现了 SSE 流或 MCP Tasks。

\section{数据格式例}
小型标准 JSON 的结构如下。示例数值只是展示维度，不用于任何物质量计算。
\par\noindent\begin{minipage}{\linewidth}
\begin{lstlisting}
{
  "schema_version": "ms-1.0",
  "time_s": [0, 1, 2],
  "mz": [44, 45],
  "signal": [[10, 9, 8], [0.12, 0.11, 0.10]],
  "signal_unit": "a.u.",
  "time_origin": "declared common experimental origin",
  "instrument_id": "example-instrument",
  "conditions_id": "example-condition",
  "synthetic": true
}
\end{lstlisting}
\end{minipage}\par
CSV 使用唯一的 \code{time_s}、\code{mz_44} 等列名，可附 \code{current_A}。CSV 导入所需单位和时间原点由调用参数提供；复杂元数据、非同网格 EC 和掩码宜使用标准 JSON。原始厂商文件应先由经验证的适配器转换，本版没有声称自动读取所有仪器格式。

\section{报告与配图构建}
\par\noindent\begin{minipage}{\linewidth}
\begin{lstlisting}[language=bash]
# Exact Arial required by default
python scripts/ms_simulate.py
# Explicit fallback only when Arial is absent
python scripts/ms_simulate.py --allow-font-fallback
OPENBLAS_NUM_THREADS=1 OMP_NUM_THREADS=1 \
  python benchmarks/ms_benchmark.py
python scripts/report_metrics.py
cd docs/latex
latexmk -xelatex -interaction=nonstopmode main.tex
\end{lstlisting}
\end{minipage}\par
若需同时重跑参考 agent，可给模拟脚本增加 \code{--agent-workspace} 参数；该目录必须没有其它服务占用。基准结果存在时才生成加速与内存图，因此重测基准后可再运行模拟脚本更新对应图。

图表默认配置为 Arial，严格模式不允许把替代字体标成 Arial。本次预览采用显式许可的 Arimo 回退，\code{font_receipt.json} 记录请求字体、实际字体和回退状态。源文件不嵌入可安装字体；PDF 内的必要渲染字形不是独立字体交付文件。

\section{复现记录的使用}
\code{VALIDATION.json} 汇总软件版本、测试计数、跳过原因、报告页数与构建状态；\code{MANIFEST.sha256} 用于核对发布文件。数值基准原始样本、模拟真值、完整恢复数组和参考调用轨迹都随包提供。重新运行可能因库版本和硬件而有小幅数值与时间差异，因此应同时记录环境，而不只比较一个加速倍数。

本版不包含前端设计章节。为保持上一版调用兼容，通用内核与原有本地界面资产仍在代码包中；质谱示例、MCP 工具目录与本报告均不再以二维或三维图像恢复为任务。

\backmatter
\bibliographystyle{unsrtnat}\bibliography{references}
\end{document}
